\section{Discussion}
\label{sec:discussion}

\paragraph{What the protocol buys.} The design choices trade statistical
power for credibility.
\begin{itemize}
  \item \emph{Counting every trial.} Invalid proposals, evaluation failures
    and degenerate trials enter the family size, and so does unused budget.
    This penalizes proposers that generate many bad candidates or stop
    early, which is exactly the behaviour that inflates apparent discovery
    rates in automated search.
  \item \emph{Counting each hypothesis once.} Trials with identical per-date
    ranks are one hypothesis. Syntactic deduplication alone would let a
    proposer inflate its discoveries by re-spelling a good expression.
  \item \emph{Confirming on unseen data.} A proposer that learns from
    formation feedback makes formation p-values optimistic; the null markets
    of Section~\ref{sec:results-synthetic} show it for genetic programming.
    Only the validation step, on data no proposer has seen, carries
    false-discovery control.
  \item \emph{The information barrier.} The harness passes proposers
    formation statistics only, and on store data later windows are not
    loaded until their phase. For an LLM, whose only input is the rendered
    prompt, this means the model cannot condition on validation or test
    outcomes of \emph{this} study. It does not remove what the model may
    have learned about the period in pre-training; that is what the
    contamination diagnostics address.
  \item \emph{The commit-then-reveal hold-out.} The ledger and the study
    registry give tamper evidence for one's own records: every commitment,
    reveal and refusal is chained, and repeated reveals are counted. They do
    not by themselves prove to a third party that no earlier reveal or
    discarded run exists, because the same process writes them. That proof
    needs the commitment hash published outside the repository, for example
    in a pushed signed tag or a public timestamp, before the separate reveal
    step, which the command-line interface enforces for real data.
\end{itemize}
The controls are designed to apply equally to every proposer, but they are
not neutral in effect: equal trial budgets are not equal compute, which is
why the design includes a compute-matched baseline arm. They make results
harder to obtain and easier to believe.

\paragraph{Reading the synthetic results.} The harness results are a
statement about the \emph{baselines} and the protocol. At a budget of
\SynBudget{} trials, undirected random search and feedback-driven genetic
programming both reach the textbook planted signals at the higher
signal-to-noise level, and both miss the library-family and the drawn
signals, which are in their search space. The selected factors are true
discoveries by the benchmark's ground truth, but mostly substantially
correlated with library factors. The benchmark thus has room to show an advantage from
better proposals, and a test of whether an advantage is more than prior
knowledge.

Such an advantage would only be informative under three conditions:
\begin{enumerate}
  \item it appears on the drawn signal, not only on the textbook and
    library-family ones that an LLM may know from the literature;
  \item it survives the same multiple-testing accounting and confirmation;
  \item it replicates across the pre-registered seeds (H2) and holds against
    compute-matched baselines.
\end{enumerate}

\paragraph{What would count as evidence for LLM-guided search.} By the
pre-registered rules, a positive answer to RQ1 requires three things:
\begin{itemize}
  \item a sealed-test improvement over both baselines at an equal budget, on
    both the market-sampling and the search-variability estimands (H1);
  \item an equivalence result showing that the improvement does not shrink
    after the model's training cutoff relative to genetic programming (H3);
  \item ideally, confirmation on a prospective window that did not exist
    when the factor set was committed.
\end{itemize}
An improvement confined to factors with low behavioural novelty would count
as rediscovery, not discovery. \todo{Discuss the actual outcome of H1 to H4
once the pre-registered runs exist.}

\paragraph{Rationales are claims, not evidence.} The LLM proposer attaches an
economic rationale and a mechanism label to every expression. These texts
are recorded but play no role in selection. Fluent but unsupported
explanations are a known failure mode of language
models~\citep{ji2023hallucination}. Whether rationales carry information
about out-of-sample survival is itself an empirical question for the
rationale ablation. \todo{Report whether rationale presence or content
predicts test survival.}

\paragraph{Scope and use.} This is research software. It places no orders
and makes no claim that any discovered factor would remain profitable after
costs, capacity limits or publication. Market prices in the real-data study
come only from a commercial data provider (MarketData) through the suite's
gateway, and no vendor data are redistributed.
